\scanpage{117}

% manuscript lines / position: 1–2
et \deleted{\unsure{les|ici}} on a une \deleted{\unsure{bonne}} op. de $\unsure{\gamma|\mathcal G|\Gamma}$ sur $A$, telle que

% manuscript lines / position: 2
\[
 \bigotimes A_i\unsure{\longrightarrow|\mathord{-}}A
\]

% manuscript lines / position: 3–5
soit compatible avec des opérations de $\gamma$, savoir celle qui fait opérer \uline{$\tau_i$} par $\sigma_{A/A_i}$.

% manuscript lines / position: 6–10
\deleted{\unsure{Considér|Considérons}} Supposons maintenant que les $A_i$ correspondent \deleted{\unsure{à un|à}} des rev. quadratiques $X_i$ de $X$, donc on a un rev. quadratique $X_1\wedge X_2\wedge X_3$, d'algèbre $B=A_1\unsure{*|\star|\wedge}A_2*A_3$,

% manuscript lines / position: 11–12
\[
 \begin{aligned}
 \text{et on a}\qquad X_1\times X_2\times X_3&\longrightarrow X_1\wedge X_2\wedge X_3,\\
 (x_1,x_2,x_3)&\longrightarrow x_1*x_2*x_3,
 \end{aligned}
\]

% manuscript lines / position: 13
\[
 \text{i.e.}\qquad B\longrightarrow A_1\otimes A_2\otimes A_3,
\]

% manuscript lines / position: 14
et de façon plus précise, on aura

% manuscript lines / position: 15–16
\[
 B\unsure{\longrightarrow|\mathord{-}}(A_1\otimes A_2\otimes A_3)^\gamma
 \qquad\substack{\text{invariants}\\\text{sous }\gamma}
\]

% manuscript lines / position: 17–20
Et \unsure{$=$|de même|même}, si on \unsure{exclut|a exclut} les $x\in X=\operatorname{Spec}(k)$ de \uline{car. 2} tels que deux au moins des $A_i$ soient non étales en ce point, alors l'hom précédent est un \uline{iso}.

% manuscript lines / position: 21–22
On va \uline{regarder} le composé

% manuscript lines / position: 22
\[
 B\unsure{\mathord{-}|\longrightarrow}A_1\otimes A_{\unsure{2|3}}\otimes A_3\unsure{\mathord{-}|\longrightarrow}A,\qquad\deleted{\unsure{\text{i.e.}}}
\]

% manuscript lines / position: bottom left
% archive sheet number: 116
